% --- Archon physics formalization source begin ---
% archon:physics
% archon:covers PhyXMiniProblems/problem_phyx_mini_0792.lean
% archon:source-report reports/phyx_mini/problem_phyx_mini_0792.source.json

\chapter{Physics problem phyx\_mini\_0792}
\label{ch:PhyXMiniProblems_problem_phyx_mini_0792}

\paragraph{Problem source.}
\#\# Physical scenario

Vector \textbackslash{}( \textbackslash{}vec\{A\} \textbackslash{}) has magnitude 6 units and is in the direction of the \textbackslash{}( +x \textbackslash{})-axis. Vector \textbackslash{}( \textbackslash{}vec\{B\} \textbackslash{}) has magnitude 4 units and lies in the xy-plane, making an angle of \textbackslash{}( 30\textasciicircum{}\textbackslash{}circ \textbackslash{}) with the \textbackslash{}( +x \textbackslash{})-axis as shown in figure.

\#\# Question

Find the vector product \textbackslash{}( \textbackslash{}vec\{C\} = \textbackslash{}vec\{A\} \textbackslash{}times \textbackslash{}vec\{B\} \textbackslash{}).

\#\# Answer choices

- A: A: \textbackslash{}[15\textbackslash{}hat\{k\}\textbackslash{}]

- B: B: \textbackslash{}[12\textbackslash{}hat\{k\} \textbackslash{}]

- C: C: \textbackslash{}[ 8\textbackslash{}hat\{k\} \textbackslash{}]

- D: D: \textbackslash{}[18\textbackslash{}hat\{k\} \textbackslash{}]

\#\# Figure caption (auxiliary; use the image as primary evidence)

The image depicts a three-dimensional coordinate system with vectors and a plane. Here's a detailed description:

1. **Axes**:
   - The coordinate system axes are labeled as \textbackslash{}(x\textbackslash{}), \textbackslash{}(y\textbackslash{}), and \textbackslash{}(z\textbackslash{}).
   - The \textbackslash{}(x\textbackslash{})-axis extends to the right, the \textbackslash{}(y\textbackslash{})-axis extends upward and out of the screen, and the \textbackslash{}(z\textbackslash{})-axis extends diagonally downward to the left.

2. **Vectors**:
   - **Vector \textbackslash{}(\textbackslash{}vec\{A\}\textbackslash{})**: Points along the \textbackslash{}(x\textbackslash{})-axis starting from the origin \textbackslash{}(O\textbackslash{}) and is represented in brown.
   - **Vector \textbackslash{}(\textbackslash{}vec\{B\}\textbackslash{})**: Extends from the origin \textbackslash{}(O\textbackslash{}) and makes an angle with the \textbackslash{}(x\textbackslash{})-axis within the plane shown. It’s also represented in brown.
   - **Vector \textbackslash{}(\textbackslash{}vec\{C\}\textbackslash{})**: Points diagonally from the origin \textbackslash{}(O\textbackslash{}) into the \textbackslash{}(z\textbackslash{})-axis direction.

3. **Plane**:
   - A beige-colored plane is shown intersecting the \textbackslash{}(x\textbackslash{}) and \textbackslash{}(y\textbackslash{}) axes. It appears to be a vertical plane containing vectors \textbackslash{}(\textbackslash{}vec\{B\}\textbackslash{}) and \textbackslash{}(\textbackslash{}vec\{A\}\textbackslash{}).

4. **Angle**:
   - An angle \textbackslash{}(\textbackslash{}phi = 30\textasciicircum{}\textbackslash{}circ\textbackslash{}) is marked between vectors \textbackslash{}(\textbackslash{}vec\{A\}\textbackslash{}) and \textbackslash{}(\textbackslash{}vec\{B\}\textbackslash{}) within the beige plane. 

5. **Origin**:
   - The point of intersection of all the vectors and axes is labeled \textbackslash{}(O\textbackslash{}).

This diagram likely represents a vector projection or interaction in three-dimensional space, with specific interest in the angle between vectors \textbackslash{}(\textbackslash{}vec\{A\}\textbackslash{}) and \textbackslash{}(\textbackslash{}vec\{B\}\textbackslash{}).

\#\# Dataset metadata

Dynamics; Spatial Relation Reasoning, Physical Model Grounding Reasoning

\paragraph{Recorded answer/context.}
B: B: \textbackslash{}[12\textbackslash{}hat\{k\} \textbackslash{}]

\paragraph{Figure/image path.}
/root/proposal\_for\_physic/science-mango/phyx\_mini\_run/phyx\_data/test\_image/792.png

\paragraph{Formalization target.}
create a compiling Lean file with sorry bodies at `PhyXMiniProblems/problem\_phyx\_mini\_0792.lean`. The Lean declarations must preserve the physical quantities, dimensions or dimensional roles, figure labels, governing-law hypotheses, and final relation expressed by this problem.
Use Mathlib/Physlib names found through LeanExplore where available. If a physics API is missing, introduce faithful local abstractions rather than scalar placeholder aliases.

\begin{theorem}[Physics formalization target]
\label{thm:physics:phyx_mini_0792:target}
\lean{PhyXMiniProblems.ProblemPhyXMini0792.vectorProduct_eq_twelve_kHat}
\uses{def:physics:phyx-mini-0792:phyxminiproblems-problemphyxmini0792-khat, def:physics:phyx-mini-0792:phyxminiproblems-problemphyxmini0792-vectorproductsetup, def:physics:phyx-mini-0792:phyxminiproblems-problemphyxmini0792-statedvectordata, def:physics:phyx-mini-0792:phyxminiproblems-problemphyxmini0792-matchessuppliedfigure, def:physics:phyx-mini-0792:phyxminiproblems-problemphyxmini0792-satisfiesvectorproductlaw}
If \(A\) has magnitude \(6\) along \(+x\), \(B\) has magnitude \(4\) in the
\(xy\)-plane at \(30^\circ\) counterclockwise from \(A\), and
\(C=A\times B\), then \(C=12\widehat{k}\).
\end{theorem}
\begin{proof}
The positive-axis and magnitude hypotheses give \(A=6\widehat{i}\).  The
figure converts the marked angle to \(\pi/6\), and the positive \(y\)-component
chooses the counterclockwise planar orientation.  Thus
\[
 B=4\cos(\pi/6)\widehat{i}+4\sin(\pi/6)\widehat{j}
   =2\sqrt3\,\widehat{i}+2\widehat{j}.
\]
The three-dimensional cross-product formula now gives
\[
 A\times B
 =(6,0,0)\times(2\sqrt3,2,0)=(0,0,12)
 =12\widehat{k}.
\]
The governing vector-product law identifies this vector with \(C\).
\end{proof}

% --- Archon declaration topology begin ---
\section{Declaration topology}
\begin{definition}[Spatial Vector]
  \label{def:physics:phyx-mini-0792:phyxminiproblems-problemphyxmini0792-spatialvector}
  \lean{PhyXMiniProblems.ProblemPhyXMini0792.SpatialVector}
  Numerical Cartesian components of a three-dimensional spatial vector
\end{definition}
\begin{proof}
  This is the declared model definition of Spatial Vector.
\end{proof}
\begin{definition}[Coordinate Axis]
  \label{def:physics:phyx-mini-0792:phyxminiproblems-problemphyxmini0792-coordinateaxis}
  \lean{PhyXMiniProblems.ProblemPhyXMini0792.CoordinateAxis}
  The three labelled coordinate axes in the supplied raster
\end{definition}
\begin{proof}
  This is the declared model definition of Coordinate Axis.
\end{proof}
\begin{definition}[axis Index]
  \label{def:physics:phyx-mini-0792:phyxminiproblems-problemphyxmini0792-axisindex}
  \lean{PhyXMiniProblems.ProblemPhyXMini0792.axisIndex}
  \uses{def:physics:phyx-mini-0792:phyxminiproblems-problemphyxmini0792-coordinateaxis}
  Coordinate index belonging to a labelled axis
\end{definition}
\begin{proof}
  This is the declared model definition of axis Index.
\end{proof}
\begin{definition}[axis Vector]
  \label{def:physics:phyx-mini-0792:phyxminiproblems-problemphyxmini0792-axisvector}
  \lean{PhyXMiniProblems.ProblemPhyXMini0792.axisVector}
  \uses{def:physics:phyx-mini-0792:phyxminiproblems-problemphyxmini0792-spatialvector, def:physics:phyx-mini-0792:phyxminiproblems-problemphyxmini0792-coordinateaxis, def:physics:phyx-mini-0792:phyxminiproblems-problemphyxmini0792-axisindex}
  Unit vector along a labelled positive coordinate axis
\end{definition}
\begin{proof}
  This is the declared model definition of axis Vector.
\end{proof}
\begin{definition}[i Hat]
  \label{def:physics:phyx-mini-0792:phyxminiproblems-problemphyxmini0792-ihat}
  \lean{PhyXMiniProblems.ProblemPhyXMini0792.iHat}
  \uses{def:physics:phyx-mini-0792:phyxminiproblems-problemphyxmini0792-spatialvector, def:physics:phyx-mini-0792:phyxminiproblems-problemphyxmini0792-axisvector}
  Unit vector i-hat along +x
\end{definition}
\begin{proof}
  This is the declared model definition of i Hat.
\end{proof}
\begin{definition}[j Hat]
  \label{def:physics:phyx-mini-0792:phyxminiproblems-problemphyxmini0792-jhat}
  \lean{PhyXMiniProblems.ProblemPhyXMini0792.jHat}
  \uses{def:physics:phyx-mini-0792:phyxminiproblems-problemphyxmini0792-spatialvector, def:physics:phyx-mini-0792:phyxminiproblems-problemphyxmini0792-axisvector}
  Unit vector j-hat along +y
\end{definition}
\begin{proof}
  This is the declared model definition of j Hat.
\end{proof}
\begin{definition}[k Hat]
  \label{def:physics:phyx-mini-0792:phyxminiproblems-problemphyxmini0792-khat}
  \lean{PhyXMiniProblems.ProblemPhyXMini0792.kHat}
  \uses{def:physics:phyx-mini-0792:phyxminiproblems-problemphyxmini0792-spatialvector, def:physics:phyx-mini-0792:phyxminiproblems-problemphyxmini0792-axisvector}
  Unit vector k-hat along +z
\end{definition}
\begin{proof}
  This is the declared model definition of k Hat.
\end{proof}
\begin{definition}[Coordinate Plane]
  \label{def:physics:phyx-mini-0792:phyxminiproblems-problemphyxmini0792-coordinateplane}
  \lean{PhyXMiniProblems.ProblemPhyXMini0792.CoordinatePlane}
  The three coordinate planes determined by the labelled axes
\end{definition}
\begin{proof}
  This is the declared model definition of Coordinate Plane.
\end{proof}
\begin{definition}[Lies In Coordinate Plane]
  \label{def:physics:phyx-mini-0792:phyxminiproblems-problemphyxmini0792-liesincoordinateplane}
  \lean{PhyXMiniProblems.ProblemPhyXMini0792.LiesInCoordinatePlane}
  \uses{def:physics:phyx-mini-0792:phyxminiproblems-problemphyxmini0792-spatialvector, def:physics:phyx-mini-0792:phyxminiproblems-problemphyxmini0792-axisindex, def:physics:phyx-mini-0792:phyxminiproblems-problemphyxmini0792-coordinateplane}
  A vector has zero component normal to the indicated coordinate plane
\end{definition}
\begin{proof}
  This is the declared model definition of Lies In Coordinate Plane.
\end{proof}
\begin{definition}[Points Along Positive Axis]
  \label{def:physics:phyx-mini-0792:phyxminiproblems-problemphyxmini0792-pointsalongpositiveaxis}
  \lean{PhyXMiniProblems.ProblemPhyXMini0792.PointsAlongPositiveAxis}
  \uses{def:physics:phyx-mini-0792:phyxminiproblems-problemphyxmini0792-spatialvector, def:physics:phyx-mini-0792:phyxminiproblems-problemphyxmini0792-coordinateaxis, def:physics:phyx-mini-0792:phyxminiproblems-problemphyxmini0792-axisvector}
  A nonzero vector is a positive scalar multiple of a labelled axis
\end{definition}
\begin{proof}
  This is the declared model definition of Points Along Positive Axis.
\end{proof}
\begin{definition}[Figure Vector Label]
  \label{def:physics:phyx-mini-0792:phyxminiproblems-problemphyxmini0792-figurevectorlabel}
  \lean{PhyXMiniProblems.ProblemPhyXMini0792.FigureVectorLabel}
  Vector labels printed beside the three arrows in image 792.png
\end{definition}
\begin{proof}
  This is the declared model definition of Figure Vector Label.
\end{proof}
\begin{definition}[Axis Sense]
  \label{def:physics:phyx-mini-0792:phyxminiproblems-problemphyxmini0792-axissense}
  \lean{PhyXMiniProblems.ProblemPhyXMini0792.AxisSense}
  Positive or negative sense along a labelled coordinate axis
\end{definition}
\begin{proof}
  This is the declared model definition of Axis Sense.
\end{proof}
\begin{definition}[Vector Product Figure]
  \label{def:physics:phyx-mini-0792:phyxminiproblems-problemphyxmini0792-vectorproductfigure}
  \lean{PhyXMiniProblems.ProblemPhyXMini0792.VectorProductFigure}
  \uses{def:physics:phyx-mini-0792:phyxminiproblems-problemphyxmini0792-coordinateaxis, def:physics:phyx-mini-0792:phyxminiproblems-problemphyxmini0792-coordinateplane, def:physics:phyx-mini-0792:phyxminiproblems-problemphyxmini0792-figurevectorlabel, def:physics:phyx-mini-0792:phyxminiproblems-problemphyxmini0792-axissense}
  Literal and qualitative information visible in the supplied raster
\end{definition}
\begin{proof}
  This is the declared model definition of Vector Product Figure.
\end{proof}
\begin{definition}[Vector Product Setup]
  \label{def:physics:phyx-mini-0792:phyxminiproblems-problemphyxmini0792-vectorproductsetup}
  \lean{PhyXMiniProblems.ProblemPhyXMini0792.VectorProductSetup}
  \uses{def:physics:phyx-mini-0792:phyxminiproblems-problemphyxmini0792-spatialvector, def:physics:phyx-mini-0792:phyxminiproblems-problemphyxmini0792-vectorproductfigure}
  The three vectors are independent data. In particular, vectorCInSquaredUnits is not defined to be either a cross product or an answer choice; the governing cross-product relation is an explicit law below
\end{definition}
\begin{proof}
  This is the declared model definition of Vector Product Setup.
\end{proof}
\begin{definition}[Stated Vector Data]
  \label{def:physics:phyx-mini-0792:phyxminiproblems-problemphyxmini0792-statedvectordata}
  \lean{PhyXMiniProblems.ProblemPhyXMini0792.StatedVectorData}
  \uses{def:physics:phyx-mini-0792:phyxminiproblems-problemphyxmini0792-liesincoordinateplane, def:physics:phyx-mini-0792:phyxminiproblems-problemphyxmini0792-pointsalongpositiveaxis, def:physics:phyx-mini-0792:phyxminiproblems-problemphyxmini0792-vectorproductsetup}
  Quantitative and geometric data stated in the problem prose
\end{definition}
\begin{proof}
  This is the declared model definition of Stated Vector Data.
\end{proof}
\begin{definition}[Matches Supplied Figure]
  \label{def:physics:phyx-mini-0792:phyxminiproblems-problemphyxmini0792-matchessuppliedfigure}
  \lean{PhyXMiniProblems.ProblemPhyXMini0792.MatchesSuppliedFigure}
  \uses{def:physics:phyx-mini-0792:phyxminiproblems-problemphyxmini0792-axisindex, def:physics:phyx-mini-0792:phyxminiproblems-problemphyxmini0792-vectorproductsetup}
  Image-derived facts. The positive y component disambiguates the two planar vectors making the same undirected angle with +x. No calibrated numerical fact about the drawn C arrow is asserted
\end{definition}
\begin{proof}
  This is the declared model definition of Matches Supplied Figure.
\end{proof}
\begin{definition}[Satisfies Vector Product Law]
  \label{def:physics:phyx-mini-0792:phyxminiproblems-problemphyxmini0792-satisfiesvectorproductlaw}
  \lean{PhyXMiniProblems.ProblemPhyXMini0792.SatisfiesVectorProductLaw}
  \uses{def:physics:phyx-mini-0792:phyxminiproblems-problemphyxmini0792-vectorproductsetup}
  Governing relation specified by C = A × B
\end{definition}
\begin{proof}
  This is the declared model definition of Satisfies Vector Product Law.
\end{proof}
\begin{definition}[Answer Choice]
  \label{def:physics:phyx-mini-0792:phyxminiproblems-problemphyxmini0792-answerchoice}
  \lean{PhyXMiniProblems.ProblemPhyXMini0792.AnswerChoice}
  Labels of the four displayed answer choices
\end{definition}
\begin{proof}
  This is the declared model definition of Answer Choice.
\end{proof}
\begin{definition}[displayed Coefficient]
  \label{def:physics:phyx-mini-0792:phyxminiproblems-problemphyxmini0792-displayedcoefficient}
  \lean{PhyXMiniProblems.ProblemPhyXMini0792.displayedCoefficient}
  \uses{def:physics:phyx-mini-0792:phyxminiproblems-problemphyxmini0792-answerchoice}
  Numerical coefficient of k-hat printed in an answer choice
\end{definition}
\begin{proof}
  This is the declared model definition of displayed Coefficient.
\end{proof}
\begin{definition}[displayed Answer Vector]
  \label{def:physics:phyx-mini-0792:phyxminiproblems-problemphyxmini0792-displayedanswervector}
  \lean{PhyXMiniProblems.ProblemPhyXMini0792.displayedAnswerVector}
  \uses{def:physics:phyx-mini-0792:phyxminiproblems-problemphyxmini0792-spatialvector, def:physics:phyx-mini-0792:phyxminiproblems-problemphyxmini0792-khat, def:physics:phyx-mini-0792:phyxminiproblems-problemphyxmini0792-answerchoice, def:physics:phyx-mini-0792:phyxminiproblems-problemphyxmini0792-displayedcoefficient}
  Vector printed by an answer choice, in squared vector units
\end{definition}
\begin{proof}
  This is the declared model definition of displayed Answer Vector.
\end{proof}
% --- Archon declaration topology end ---
% --- Archon physics formalization source end ---
